\clearpage\begingroup\SourceLeaf{76}\HistoricalHead{71}{R}{BERECHNUNG MAGNETISCHER KRÄFTE.}
\begin{center}[12.]\end{center}
\noindent\rule{\textwidth}{.4pt}\par
\begin{center}\small [Aus Handbuch 21, Bg. Aufsätze, Notizen und Rechnungen zur Mathematik und Astronomie gehörig.\\
Angefangen September 1813, S. 19.]\end{center}
\noindent\rule{\textwidth}{.4pt}
\begin{center}\textbf{Theorem über die Anziehung.}\end{center}
\indent Ein Körper von irgend einer beliebigen Gestalt und Zusammensetzung
übt in jedem Punkte des Raum’s eine bestimmte Anziehung aus. Sobald
man überall deren \textls[120]{Richtung} angeben kann, ist man auch im Stande ihre
Intensität anzugeben, ohne die Gestalt des Körpers selbst zu kennen. Denkt
man sich nemlich eine unendlich enge Röhre deren Wände eine solche Gestalt
haben, dass sie überall von der Richtung der Anziehung tangiert werden,
so ist die Anziehung überall der Weite proportional.

\raisebox{-.6ex}{*}NB. Der Satz ist nicht allgemein zu verstehen; bloss die Intensitäten
in \textls[120]{derselben} Röhre kann man so unter sich vergleichen, nicht aber in verschiedenen
Röhren.\textsuperscript{*}

\raisebox{-.6ex}{**}Die Vergleichung wird vollständig, indem man zugleich die Systeme
der Flächen betrachtet, gegen welche die Richtung normal ist.\textsuperscript{**}[\textsuperscript{*})]
\begingroup\renewcommand{\thefootnote}{}\footnotetext{[\textsuperscript{*}) In der Handschrift ist die oben zwischen \raisebox{-.6ex}{*}\ \textsuperscript{*} gesetzte Stelle mit anderer Tinte, die zwischen
\raisebox{-.6ex}{**}\ \textsuperscript{**} gesetzte Stelle mit Bleistift geschrieben.]}\endgroup
\bigskip\begin{center}\textbf{BEMERKUNG.}\end{center}
{\small\noindent\hspace*{\parindent}In der vorstehenden Aufzeichnung [12.] hat \textsc{Gauss} den Begriff der \textls[120]{Kraftröhre} eingeführt, sowie
das Feld \textls[120]{durch Flächen} gleichen Potentials und Kraftröhren in Zellen eingeteilt, genau so, wie es
später durch \textsc{Faraday} und \textsc{Maxwell} geläufig geworden ist. Man vergl. etwa \textsc{Maxwells} Abhandlung
\textit{On Faraday’s Lines of Force}, Transactions of the Cambridge Philos. Society \textbf{10}, 1856, S. 27, Scientific
Papers I, S. 155, deutsch in \textsc{Ostwalds} \textit{Klassikern} Nr. 69, 1895.
\par\raggedleft\textsc{Schaefer}.\par}
\bigskip\begin{center}\rule{.23\textwidth}{.4pt}\end{center}
\clearpage\endgroup
